\DTMsavedate{mydate}{2022-06-30}
\maketitle{Control Engineering}{Electronic Engineering}{\DTMusedate{mydate}}

% ----------------------------------------------------- %
% COURSE INFO
\subsection*{Course Information}\label{course:3241}
\vspace{0ex}%
\noindent\parbox{0.5\textwidth}{%
    \noindent\begin{tabular}{@{}ll}
                 \textsf{Course:}          & 	Control Engineering \\
                 \textsf{Course Code:}         & 	TEE 3241    \\
                 \textsf{Credit Hours:}    & 	10
    \end{tabular}}
\vspace{2ex}\hrule\vspace{2ex}

% ----------------------------------------------------- %
% INSTRUCTOR INFO
\subsection*{Lecturer Information}
\noindent\parbox{0.5\textwidth}{%
    \noindent\begin{tabular}{@{}ll}
                 \textsf{Name:}         &     Reginald Gonye         \\
%\textsf{Phone:}        &    \phone          \\
\textsf{Email:}        &    reginald.gonye@nust.ac.zw \\
\textsf{Qualifications:}        &  MPhil in Electronic Engineering, (Hon) BEng in Electronic Engineering
    \end{tabular}}
\vspace{2ex}\hrule\vspace{2ex}


% ----------------------------------------------------- %
% COURSE SYNOPSIS
\subsection*{Course Synopsis}
The module explores examples of controlled processes, objectives and terminology, open and closed-loop controllers.
It covers modeling by transfer functions, simple servo mechanisms and derivation of transfer functions from specifications.
It also covers time and frequency–response specifications, direct analysis and design, stability, Routh criterion, the ITAE and other performance criteria.
There are examples of servo design,  frequency-response analysis and design, Root-locus methods and system analysis and design.
\vspace{2ex}\hrule\vspace{2ex}


% ----------------------------------------------------- %
% COURSE DESCRIPTION
\subsection*{Learning Aim}
To instill concepts of Control Engineering Theory into students.
\vspace{2ex} \hrule\vspace{2ex}

% ----------------------------------------------------- %
% LEARNING OBJECTIVES
\subsection*{Learning Objectives}
By the end of the course, candidates should be able to:
\begin{outline}
    \1      Understand the concept of a control system.

    \1  	Understand modelling of control systems.

    \1  	Use system models to analyse and synthesise systems.

\end{outline}
\vspace{2ex}\hrule\vspace{2ex}

% ----------------------------------------------------- %
% COURSE TOPICS
\subsection*{Topics}
\begin{outline}
    \1 Introduction
    \1 Modeling Control Systems
    \1 Time Response of a System
    \1 Performance Criteria
    \1 Controllers and Compensators
    \1 Stability of a System
    \1 The Root Locus
    \1 Frequency Response
    \1 Introduction to Digital Control
    \1 Application of PLCs in Control Systems
    \1 Case studies of Control Systems
\end{outline}
\vspace{2ex} \hrule\vspace{2ex}

% ----------------------------------------------------- %
% Students Assenssment
\subsection*{Assessment of Students}
\begin{outline}
    \1 Continuous assessment will consist of tests and assignments (NB: Attendance may also contribute to the coursework mark).
    \1 Final assessment will consist of an examination at the end of the semester. The examination shall comprise of two sections: section A carrying 40 marks and section B with five questions carrying 20 marks each. You will be required to answer all questions in section A and to attempt any three questions from section B. Section A questions will test factual recall of control concepts in form of calculations or accounts. Section B questions will test your ability to apply the concepts learnt in the course to control systems.
\end{outline}
\vspace{2ex}\hrule\vspace{2ex}

% ----------------------------------------------------- %
% Grade Evaluation
\subsection*{Course Evaluation and Grading Policies}
Grades are based on:
Continuous assessment (\qty{25}{\percent}),
Final Exam (\qty{75}{\percent})
\vspace{2ex}\hrule\vspace{2ex}
% ----------------------------------------------------- %
% EQUIPMENT
% https://sites.udel.edu/computing-purchases/personal-specs/
% \subsection*{Essential Equipment}

% \vspace{2ex}\hrule\vspace{2ex}

% Policies and Other information
\subsection*{Policies and Other Information}
% Key Text or some Books
\subsection*{Key Text}

\begin{itemize}
    \item Roland S Burns, 2001, Advanced Control Engineering, Butterworth-Heinemann
    \item D K Anand and R B Zmood, 1995, Introduction to Control Systems Third Edition, Butterworth-Heinemann
\end{itemize}
\vspace{2ex}\hrule
\vspace{2ex}\hrule\vspace{2ex}